\chapter{Introduction and Approach}
\label{ch:supp:intro}

% Framing chapter covering scope, evidence base, the PUE definition and its limits,
% and the standard of inference adopted (logic qualified per review).

This document provides Protocol Policy Lab's answers to the two supplementary questions issued following the hearing of 1 May 2026. It is lodged as a standalone response. The inquiry to which it contributes is itself a novel exercise, described by its proponents as a process to \enquote{establish a parliamentary inquiry into data centre developments in NSW, the first of its kind in Australia} \autocite{greensNSWDataCentreInquiry2026}. That novelty matters for both questions, because it means New South Wales is forming its evidentiary and regulatory baseline at the same moment that the sector is expanding fastest.

\section{Scope and the questions answered}
\label{sec:supp:scope}

The first supplementary question asks whether Protocol holds data on the typical Power Usage Effectiveness (PUE) of data centres currently operating in New South Wales, whether such information is publicly available, and whether Protocol supports mandatory PUE reporting as a baseline transparency measure. The second asks whether Protocol is aware of evidence or analysis bearing on whether energy gentrification is occurring or likely to occur in New South Wales, with particular reference to Western Sydney, where data centre clustering is most concentrated. \Cref{ch:supp:pue} answers the first and \Cref{ch:supp:gentrification} the second. \Cref{ch:supp:synthesis} draws out the single thread that connects them, which is that neither efficiency policy nor consumer protection can operate ahead of measurement and disclosure.

\section{Evidence base and method}
\label{sec:supp:method}

The answers draw on two bodies of evidence. The first is the peer-reviewed and technical literature on data centre energy, water and grid integration. The second is the primary and grey literature specific to the Australian National Electricity Market and to New South Wales, including the Australian Energy Market Operator forecasts as relayed in the New South Wales Government's own consultation paper, transmission planning by Transgrid, and analysis by Energy Consumers Australia and the United States Studies Centre. Quotations throughout are reproduced verbatim from the cited source. Measured quantities are distinguished from projections, and projections are attributed to their modelling scenario wherever the source specifies one. Where no New South Wales figure exists, this is stated rather than inferred from overseas data.

\section{Power Usage Effectiveness and the limits of a single metric}
\label{sec:supp:pue-defn}

Power Usage Effectiveness is the ratio of a facility's total energy to the energy delivered to its information technology equipment. A value of 1.0 is the theoretical ideal in which no energy is spent on overheads, and lower values denote greater efficiency. The metric is mature and standardised, having been established such that \enquote{the assessment method of Power Usage Effectiveness has become an international standard for measuring energy efficiency and is widely used in the design and operation of data centers} \autocite{uanovAnalysisMethodsAssessing2022}. Cooling is the dominant overhead the ratio captures, since \enquote{the energy consumption of the cooling system for the machine room accounts for a significant part of the operating costs of the building} \autocite{uanovAnalysisMethodsAssessing2022}.

PUE is a necessary instrument but an insufficient one, and this submission is careful not to treat it as a complete measure. It is a ratio, so it is silent on absolute consumption, on the carbon intensity of the electricity drawn, and on water. The water dimension is consequential and under-measured, given that \enquote{Data centres consume water directly for cooling, in some cases 57\% sourced from potable water} \autocite{myttonDataCentreWater2021}, and that the academic literature records how, for data centres, \enquote{their water footprint (WF) has so far received little to no attention} \autocite{risticWaterFootprintData2015}. Water intensity is also enormously site-dependent, with one assessment finding \enquote{workload-level water use variations exceeding 10,000-fold} \autocite{leiWaterUseData2025}. A PUE figure can improve while absolute energy, emissions and water all rise, which is why the recent literature advocates \enquote{mandatory resource transparency, efficiency regulation, and lifecycle accountability} rather than a single efficiency ratio \autocite{imanHiddenCostsIntelligence2025}.

The insufficiency is compounded by rebound. Efficiency gains at the facility level are routinely offset by growth in demand, so that \enquote{the energy-intensive nature of digital infrastructure may conversely increase emissions}, with the gains from efficiency repeatedly overtaken by rising consumption \autocite{liuImpactGlobalDigital2025}. The United States experience illustrates the same point in reverse, where for two decades \enquote{comparable increases in electricity demand have been avoided through the adoption of key energy efficiency measures and a shift towards large cloud-based service providers} \autocite{shehabiDataCenterGrowth2018a}, a stabilisation that the artificial intelligence load is now ending. The practical implication for New South Wales is that an efficiency-conditioned regime should require a basket of measured indicators, being PUE alongside water usage effectiveness, absolute energy and renewable share, not PUE in isolation.

\section{The standard of inference adopted}
\label{sec:supp:inference}

Two of the conclusions in this submission rest on inference rather than direct local measurement, and the standard applied to them is stated plainly here so that the Committee can weigh them appropriately. For the question of mandatory reporting, the argument is that measurement to a common standard and disclosure are a precondition for any efficiency-conditioned or capacity-allocation policy, since a standard that is not measured cannot be enforced and a standard that is not disclosed cannot be scrutinised. For the question of energy gentrification, the argument is deliberately bounded. Protocol does not assert that energy gentrification is occurring in New South Wales, nor that it will occur. The claim is the weaker and more defensible one that the conditions which preceded the effect in comparable jurisdictions are present in Western Sydney, that the outcome is therefore a material and foreseeable risk, and that because the outcome is contingent on policy settings it is preventable. This framing is adopted precisely because the contingency strengthens rather than weakens the case for early action, since a preventable harm warrants pre-emptive safeguards while it remains preventable.
